%! Author = sriadityavedantam
%! Date = 3/30/26

\section{Simplicity of $A_n$}

\begin{theorem}{
    If $n\geq 5$, then $A_n$ is simple.
}
Outline of proof.

\textbf{1.}
$A_n$ is generated by $3$-cycles.
Indeed, every element of $A_n$ is, by definition, a product of an even number of transpositions.
Every pair of transpositions is of one of the following forms:
\[
    (ab)(cd), \qquad (ab)(ac), \qquad (ab)(ab).
\]
In the first case,
\[
    (ab)(cd)=(adb)(adc).
\]
In the second case,
\[
    (ab)(ac)=(acb).
\]
In the third case,
\[
    (ab)(ab)=e=(abc)(acb).
\]
Thus every pair of transpositions is either a $3$-cycle or a product of two $3$-cycles.
Hence every element of $A_n$ is a product of $3$-cycles.

\textbf{2.}
Suppose $N\trianglelefteq A_n$ and $(123)\in N$.
Then $N$ contains every $3$-cycle.
Indeed, since $N$ is normal, for every $x\in A_n$,
\[
    x(123)x^{-1}\in N.
\]
But conjugation preserves cycle structure, and
\[
    x(123)x^{-1}=(x(1)\ x(2)\ x(3)).
\]
Since every $3$-cycle in $A_n$ is of this form, $N$ contains all $3$-cycles.
Because $A_n$ is generated by $3$-cycles, this implies
\[
    N=A_n.
\]

\textbf{3.}
Now suppose $1\neq N\trianglelefteq A_n$.
We show that $N$ contains a $3$-cycle.
Take
\[
    1\neq \sigma\in N.
\]
Write $\sigma$ as a product of disjoint cycles.
If one of the disjoint cycles has length at least $3$, say
\[
    (1\,2\,3\,\dots\,r)
\]
with $r\geq 3$, let
\[
    \tau=(1\,2\,3).
\]
Then the commutator
\[
    \tau\sigma\tau^{-1}\sigma^{-1}\in N
\]
since $N$ is normal.
A direct computation shows that this produces a $3$-cycle.
For example, when the support is chosen appropriately, one gets an element such as
\[
    (1\,2\,4)\in N.
\]
If instead $\sigma$ is a product only of disjoint $2$-cycles, then because $\sigma\in A_n$, there are an even number of them.
Since $n\geq 5$, after relabeling we may assume
\[
    \sigma=(1\,2)(3\,4)\cdots.
\]
Let
\[
    \tau=(1\,2\,3)\in A_n.
\]
Then again the commutator
\[
    \tau\sigma\tau^{-1}\sigma^{-1}\in N
\]
is a $3$-cycle.
Thus in every case, $N$ contains a $3$-cycle.
By Step 2, it follows that
\[
    N=A_n.
\]
Therefore $A_n$ is simple for all $n\geq 5$.
\end{theorem}

\begin{theorem}{
    For each $n\geq 5$, the alternating group $A_n$ is simple.
    Also, $A_3$ is simple.
}
The case $n\geq 5$ is the previous theorem.
Also,
\[
    A_3=\{e,(123),(132)\},
\]
so
\[
    \ord(A_3)=3.
\]
Since $3$ is prime, any subgroup of $A_3$ has order $1$ or $3$ by Lagrange's Theorem.
Thus $A_3$ has no nontrivial proper normal subgroups.
Hence $A_3$ is simple.
Note that $A_1$ and $A_2$ are trivial, so they are not simple under our definition, and $A_4$ is not simple.
\end{theorem}